\section{The two family indices}
Circle exchange compares the stable-pair theories on the fixed threefolds of two different fiber rotations of a fivefold. Let $C$ be a smooth projective complex curve of genus $g=b+1\ge1$,
let $B=\Pic^b(C)$, and let $L$ be the Poincar\'e line on $C\times B$
trivialized at a chosen point of $C$. Put $K=K_C$, $M=K/L$ and
$D(F)=R\pi_*F$ for the projection $\pi:C\times B\to B$. Products
and quotients of lines denote tensor products and duals. For a
threefold whose normal bundle over $C$ is $L\oplus M$, the equivariant
stable-pair index can be compared with that of the product
$C\times\AAo^2$. The comparison varies $L$ over $B$ on both sides.

Consider a double curve in genus two. At the minimum Euler grading
the full normal torus fixes two pairs over each $L$, the ribbons in
the two normal directions. They have Euler characteristic $-3$. Write $q$ for the original Euler grading, $t$
for the exchanged grading, and $\delta=1-\kappa$ for the change in
the volume character. At $q^{-3}t^2$ the coefficient of $\delta$ is
$-6\psi^2 D(L)$, where $\psi^j$ denotes the integral Adams operation.
This class has rank zero. Its Chern character in rational singular
cohomology is $12\theta$, where $\theta$ is the principal polarization
of $\Pic^1(C)$. Indeed Grothendieck--Riemann--Roch in cohomology gives
$\operatorname{ch}D(L)=-\theta$.
Passing to the numerical index of a single target loses this coefficient.
The calculation at $q^{-3}$ raises two questions. Does the variation
agree at every Euler grading and in arbitrary curve degree? How much
of the comparison persists when $\kappa$ is free?

The two rotations act on the ordered summands of the fivefold
\[
 \operatorname{Tot}_{C\times B}
 (tL\oplus(\kappa/t)M\oplus q\OO\oplus(c/q)\OO),\qquad
 c=\kappa^{-1},\quad \sigma^2=\kappa .
\]
The original threefold uses the first two summands, the rotated one
the last two. A stable pair is a section of a pure one-dimensional
sheaf with proper support and zero-dimensional cokernel. Its degree
$m$ is the multiplicity over $C$. In the original index a pair of
Euler characteristic $j$ is counted with $(-q)^j$. In the rotated
index the corresponding monomial is $(-t)^k$, where
$k=\chi(\mathcal F)+mb$. We call $[t]$ the one-mark coefficient: on the product threefold a single box of the monomial quotient acquires a divisor of length one. The modified virtual structure sheaf uses
the ordered external lines in each presentation, as specified in
Section~\ref{sec:geometry}. Denote the resulting pushforwards to $B$
by $\Iorig,\Irot$. In a partition $\lambda$ write $a_x,\ell_x,h_x$
for the arm, leg and hook length $1+a_x+\ell_x$ of a box $x$. Set
\[
 Z_\lambda(q)=\prod_{x\in\lambda}[(1-q^{h_x})(1-q^{-h_x})]^b,\qquad
 D_m=bm(m+1)/2.
\]

\begin{theorem}\label{thm:first}
For every $m\ge1$, both $\sigma^{mb}\Iorig$ and $\sigma^{mb}\Irot$
belong to $q^{-D_m}K^0(B)[\kappa^{\pm1}][[q,t]]$ and are congruent
modulo $(1-\kappa)^2$. For $k>0$ the common first variation is
\begin{equation}\label{eq:first}
 [\delta t^k]\mathcal I_m
 =-\sum_{\lambda\vdash m}Z_\lambda(q)
 \sum_{\substack{x\in\lambda\\h_x\mid k}}
 \psi^{k/h_x}D(L^{h_x}K^{-\ell_x}).
\end{equation}
The formula holds after restriction to any orbit of diagrams under
transposition. Here the positive-$t$ coefficient is divided by
$1-\kappa$ before specializing $\sigma=1$. At $t^0$, on such an
orbit, put $z_j=[q^j]\sum_\lambda Z_\lambda(q)$. The coefficient of
$q^j$ in either unnormalized character is
$z_j(1+j(\sigma-1))$ modulo $(\sigma-1)^2$.
\end{theorem}

For the second comparison set $u=q$, $v=c/u$, $p=K$, $s=p/c$,
$R(z)=(1-sz)/(1-z)$ and $d_{ij}=u^iv^j$, with $(0,0)$ the
northwest box, $i$ increasing to the right and $j$ downwards. Denote the removable and addable boxes of $\lambda$
by $\mathcal R_\lambda$ and $\mathcal A_\lambda$. Define
\begin{align}
 Z_\lambda(q,\kappa)
 &=(-1)^{mb}q^{-b\sum_x h_x}\kappa^{-b\sum_x(\ell_x+1/2)}
 \prod_x(1-\kappa^{\ell_x}q^{h_x})^b
         (1-\kappa^{\ell_x+1}q^{h_x})^b,\label{eq:empty}\\
 G_\lambda
 &=R(u)R(v)\sum_{r\in\mathcal R_\lambda}
 \frac{\prod_{h\in\mathcal R_\lambda\setminus\{r\}}R(d_r/d_h)}
 {\prod_{a\in\mathcal A_\lambda}R(cd_r/d_a)}.\label{eq:corners}
\end{align}

\begin{theorem}\label{thm:mark}
At unrestricted volume, for every $m\ge1$, the contributions of each
$\lambda\vdash m$ to $[t]\Iorig$ and $[t]\Irot$ agree at every Euler
grading. Their common sum in $q^{-D_m}K^0(B)[\sigma^{\pm1}][[q]]$ is
\begin{equation}\label{eq:mark}
 [t]\mathcal I_m
 =\sum_{\lambda\vdash m} Z_\lambda(q,\kappa)\,
 D\bigl(L(1-c)G_\lambda\bigr).
\end{equation}
The summand is the contribution of the indicated diagram. Its
quotients are expanded at $q=0$ before pushforward.
\end{theorem}

\begin{corollary}\label{cor:strip}
For each $\lambda$ the localized family characters agree at
$q^{-b\sum_{x\in\lambda}h_x+n}t^k$ whenever $n,k\ge0$ and
$\min(n,k)\le1$.
\end{corollary}

In \cite[Theorem~1.1]{W} the two characters at $\kappa=1$ agree
already on the common space of cycles. They are supported on the
section $m[C]$. After pushforward to $B$, Theorem~\ref{thm:first} shows that the
equality persists modulo $(1-\kappa)^2$ with the common normalization. It also extends
the genus-two, degree-two calculation of \cite[Theorem~1.2]{W},
which treated the lowest Euler grading and every $k>0$.
Theorem~\ref{thm:mark} instead retains the entire volume character
at $k=1$. Both theorems here concern classes on $B$. Lifting these
refined equalities to the common space of cycles is a further question.

The two computations of \eqref{eq:first} have different fixed loci.
On the local curve the fixed virtual sheaf has dimension zero. On
the product threefold a moment-map deformation separates the boxes
into walls whose moduli are symmetric powers of $C$. For
\eqref{eq:mark} the distinction occurs already on a disk. One
expression is a Macdonald descendant of the vertex of
\cite{DS}, the other is the sum of individual modifications at the
corners. Their equality follows by telescoping the contributions of
pairs of boundary steps. Keeping both divisors of marks in the
descent gives the stated Picard classes.
